\chapter*{Declaración de uso de inteligencia artificial generativa}
\addcontentsline{toc}{chapter}{Declaración de uso de inteligencia artificial generativa}

En la elaboración de este Trabajo Fin de Máster he utilizado herramientas de
inteligencia artificial generativa como apoyo técnico, bibliográfico y
documental. Las principales herramientas empleadas han sido Claude y Claude
Code (Anthropic), y ChatGPT y Codex (OpenAI), con las siguientes finalidades:

\begin{itemize}
  \item \textbf{Desarrollo de CitaLab.} Asistencia en la elaboración, revisión
        y depuración del código, las pruebas y los programas de análisis, así
        como en la preparación y revisión de las instrucciones de los agentes.
  \item \textbf{Apoyo bibliográfico.} Asistencia para orientar búsquedas,
        localizar trabajos relacionados, elaborar resúmenes preliminares y
        comparar sus aportaciones. Estos materiales se han utilizado como ayuda
        para el estudio de la literatura, contrastando las referencias y las
        afirmaciones incorporadas con las publicaciones originales.
  \item \textbf{Documentación y memoria.} Apoyo en la elaboración, organización
        y revisión de la documentación técnica y de la memoria, incluyendo la
        estructura de los contenidos, el estilo, la coherencia y las
        referencias.
  \item \textbf{Evaluación experimental.} Los modelos utilizados en el
        asistente, el usuario simulado y el jurado forman parte del sistema y
        del método experimental descritos en la memoria. En particular, Claude
        Sonnet 5.5 y GPT-6.1-Terra constituyen el jurado que puntúa las
        conversaciones. Esta evaluación mediante modelos se identifica y se
        distingue de la revisión humana.
\end{itemize}

He mantenido la supervisión y el criterio sobre el trabajo, revisando
críticamente las propuestas y el material incorporado. Las decisiones finales
sobre los objetivos, el diseño, la metodología y la interpretación de los
resultados han sido adoptadas por mí. Asumo la responsabilidad sobre el
código, las fuentes, los resultados y el texto presentado.
